\documentclass[../main.tex]{subfiles}
\begin{document}

\section{Discussion}
\label{sec:discussion}

The robust-estimation comparison cautions against treating MIS-selected observations as a deletion recommendation. Under response-outlier and good-leverage contamination, deleting the MIS-selected set can increase coefficient error, whereas under bad leverage it substantially improves accuracy. Even in the latter design, however, post-deletion OLS remains less accurate than MM regression and LTS. More generally, when contamination is spread across the sample rather than concentrated within the deletion budget, removing the most influential set may leave other sources of estimation error unresolved. MIS therefore identifies where target-coefficient sensitivity is concentrated, but does not guarantee estimator stability or provide an automatic data-cleaning rule.

Model misspecification provides a second source of concentrated target-coefficient sensitivity. Nonlinear confounding, heterogeneous slopes, nonlinear functional form, and threshold behaviour generate strong sensitivity responses across the recorded environments, while omitted-variable misspecification is substantially weaker in the simulated design. The linear-endogeneity design illustrates a formal limit of the diagnostic. Its severity component lies within the fitted linear span, so increasing severity changes the coefficient representation while leaving the standardised deletion-sensitivity structure unchanged. The occurrence of concentrated coefficient sensitivity is therefore mechanism specific under the stated deletion budget.

The structure of the selected set can also vary independently of the strength of the sensitivity signal. In the threshold design, MIS-selected observations are substantially enriched for membership in the affected subgroup, while heterogeneous slopes produce much weaker enrichment despite substantial coefficient sensitivity. The post-deletion results further show that coefficient RMSE can increase as the deletion budget grows. Thus, locating influential observations does not by itself resolve misspecification within the maintained model. It localises where finite-sample coefficient sensitivity is concentrated, while the substantive meaning of that localisation depends on the mechanism generating the model departure.

The empirical audits show that target-coefficient sensitivity varies in both magnitude and structure across realised estimation samples. Eight of the ten selected coefficients cross zero along at least one searched deletion path within an approximately \(5\%\) deletion budget. The first zero crossing ranges from one observation, approximately \(0.009\%\) of the estimation sample in the cash-versus-in-kind application, to approximately \(3.46\%\) in the electricity-information application. Influential-set composition also varies across applications, deletion budgets, and search directions, including persistent influential cores, membership turnover across budgets, and geographic concentration. These audits complement the controlled simulations by documenting heterogeneous patterns of finite-sample coefficient dependence in realised data.

\subsection{Limitations}

Several scope conditions delimit the interpretation of these results. The deletion size \(k\) is prespecified and functions as a sensitivity budget; different values can produce different influential sets. MIS is also coefficient-specific, so the selected set can change with the target coefficient even within the same fitted model. Global optimisation in this thesis applies to the stated fixed-\(k\), fixed-residualised objective. Full-refit validation evaluates the selected candidate set after re-estimating the frozen audit specification, but it does not exhaustively optimise the refitted coefficient over all size-\(k\) deletion sets. The contamination simulations additionally use oracle-assisted deletion budgets and, in some designs, the known direction of the induced coefficient perturbation. Because coefficient movement arises across many values of $k$ and in both directions, the thesis reports deletion paths rather than selecting a single deletion budget, and it attaches no p-values to selected sets. Calibration enters only where a detection question requires it, in the misspecification study, through independently simulated null distributions; the block-maxima EVT construction examined in \autoref{app:evt} is not used for inference.

The empirical audits carry additional interpretive limits. Each result concerns one selected estimand and validated specification, so coefficient sensitivity does not by itself determine the validity of the source study's identification strategy or the sensitivity of other estimands. Several applications use documented linear companion or robustness specifications because their headline estimators lie outside the current MIS search framework. The composition of an influential set also provides descriptive evidence about where coefficient dependence is concentrated, while explanation of that influence requires study-specific information. Observable features such as outcomes, covariates, geography, institutional structure, measurement, or heterogeneous effects can help characterise the selected observations without identifying a unique causal mechanism.

\subsection{Future Development}

A first direction for future development is richer characterisation of influential-set structure within individual applications. Once MIS has located an influential set, subsequent analysis can examine baseline characteristics, treatment status, outcomes, residuals, geographic location, institutional setting, or other study-specific variables. Persistence across deletion budgets can further distinguish observations that repeatedly form an influential core from those that enter only at particular values of \(k\). Such analysis would connect coefficient sensitivity to observable empirical structure while preserving the distinction between localisation and explanation.

A second direction concerns the scope of the optimisation framework. Recent work has extended influential-set analysis from adaptive search procedures to global optimisation of fixed-\(k\) linear-regression objectives \parencite{kuschnig2021hidden,konrad2026finding}. Further work could develop coefficient-targeted subset search under structured deletion constraints and for estimands or estimators outside the current linear representation. This extension is particularly relevant to the empirical audits, where several headline specifications use instrumental-variable, GMM, or nonlinear estimators and are currently represented through documented linear companion specifications. Broader estimator coverage would allow the same coefficient-sensitivity question to be studied directly across a wider class of empirical models.

\end{document}